% =============================================================================
% Symbols and glossary
% =============================================================================
\mysection{Symbols and glossary}

\begin{description}[leftmargin=2.6cm, style=nextline]
\item[AM] Additive Manufacturing --- a process of building components layer by layer from metallic
  powder or wire feedstock using a laser or electron beam.
\item[LPBF] Laser Powder Bed Fusion --- the most common AM process, where a laser selectively melts
  successive layers of powder spread across a build platform.
\item[DED] Directed Energy Deposition --- an AM process using focused energy to melt material as it
  is deposited.
\item[EBM] Electron Beam Melting --- an AM process using an electron beam as the energy source.
\item[XCT] X-ray Computed Tomography --- a non-destructive imaging technique that reconstructs the
  three-dimensional internal structure of a component from a series of 2D X-ray projections.
\item[NIST] National Institute of Standards and Technology --- the US federal agency that provides
  the CoCr AM XCT benchmark dataset used for training and evaluation in this thesis.
\item[CNN] Convolutional Neural Network --- a deep learning architecture widely used for image
  analysis tasks, including defect detection in XCT data.
\item[U-Net] A deep learning architecture with a symmetric encoder-decoder structure and skip
  connections, used in this thesis as the primary defect detection model.
\item[NLM] Non-Local Means --- a denoising algorithm that replaces each voxel intensity with a
  weighted average of similar voxels in a search window, preserving defect boundaries while
  suppressing noise.
\item[BHC] Beam Hardening Correction --- a preprocessing step that removes the systematic intensity
  gradient caused by differential X-ray attenuation in dense metal components.
\item[SNR] Signal-to-Noise Ratio --- a quantitative measure of image quality used to evaluate the
  effectiveness of each preprocessing stage.
\item[GPU/CPU] Graphics Processing Unit / Central Processing Unit --- hardware used for training and
  running deep learning defect detection models.
\item[FAIR] Findable, Accessible, Interoperable, Reusable --- principles for data stewardship ensuring
  long-term usability and reproducibility of datasets and pipeline outputs.
\item[DVC] Data Version Control --- a tool used alongside Git to version control large XCT datasets
  and ensure reproducibility of all experiments.
\item[MLflow] An open-source platform used for experiment tracking, logging hyperparameters, metrics,
  and model checkpoints across all training runs.
\item[LoF] Lack-of-Fusion --- a type of internal defect in AM components caused by insufficient
  energy density during fabrication, resulting in irregularly shaped voids with sharp
  stress-concentrating features.
\item[Gas Porosity] A type of internal defect in AM components formed by entrapped gas in the melt
  pool, typically spherical with diameters ranging from 1 to 500 micrometres.
\item[Otsu Threshold] An automatic global thresholding method that maximises inter-class intensity
  variance. Used in this thesis as both a classical detection baseline and an automatic label
  generator.
\item[Pseudo-Label] An automatically generated binary detection label produced by classical
  thresholding (Otsu for the NIST phase, Bernsen for the PODFAM phase), used as training supervision
  in the absence of expert-annotated ground truth.
\item[Dice Coefficient] A detection accuracy metric measuring the overlap between predicted and
  reference defect regions, ranging from 0 (no overlap) to 1 (perfect overlap).
\item[IoU] Intersection over Union --- a detection accuracy metric measuring the ratio of overlap to
  union between predicted and reference defect regions.
\item[BCE] Binary Cross-Entropy --- a standard loss function for binary classification tasks. Shown
  to cause class collapse on severely imbalanced XCT defect data when used alone.
\item[Dice-Focal Loss] The combined loss function used in this thesis, combining Dice loss and Focal
  loss with equal weighting to address severe class imbalance in XCT defect detection.
\item[Patch Extraction] The process of dividing full XCT slices into fixed-size sub-images
  ($256 \times 256$ pixels in this thesis) for input to the detection model.
\item[ReLU] Rectified Linear Unit --- the activation function used in the encoder and decoder stages
  of the U-Net, introducing non-linearity at low computational cost.
\item[$\mu$m] Micrometre --- unit of length used to describe defect sizes and XCT voxel dimensions.
  $1\,\mu\text{m} = 10^{-6}$ metres.
\item[$10^6$ K/s] Cooling rate in Kelvin per second, typical of rapid solidification dynamics in LPBF
  additive manufacturing processes.
\item[Bernsen / DCT / LCT] Bernsen local adaptive thresholding, used as the primary pseudo-label
  method for the PODFAM phase (Part II). DCT (dynamic contrast threshold) and LCT (local contrast)
  are its two governing parameters, defined in Section~9.3.
\end{description}
